% -----------------------------------------------------------
\section{Resurrection}
% -----------------------------------------------------------
	\label{RESURRECTION}
	
% ----------------------
\subsection{See also}
% ----------------------

	\hyperref[DEATH]{Death},

% % ----------------------
% \subsection{1828 American Dictionary of the English Language} % *1828
% % ----------------------
	% \label{RESURRECTION-1828}

	% \begin{compactitem}
		% \item
	% \end{compactitem}

% % ----------------------
% \subsection{Oxford English Dictionary} % *OED
% % ----------------------
	% \label{RESURRECTION-OED}

	% \begin{compactitem}
		% \item[]
	% \end{compactitem}

% ----------------------
\subsection{Scriptural Usage} % *SCRIPTURES
% ----------------------
	\label{RESURRECTION-SCRIPTURES}

	% % -----
	% \subsubsection{Old Testament} % *OT
	% % -----
		% \label{RESURRECTION-OT}
	
		% % --
		% \paragraph{KJV}
		% % --
			% \label{RESURRECTION-OT-KJV}
	
			% \begin{tabular}{ll}
				% Total occurrences in the OT & 
			% \end{tabular}
			
			% \begin{compactdesc}
				% \item[]
			% \end{compactdesc}
			
		% % --
		% \paragraph{Hebrew Bible}
		% % --
			% \label{RESURRECTION-HEBREW}
		
			% %
			% \subparagraph{\texorpdfstring{\he{sample word}}{}} % paste actual hebrew text here
			% %
				% \label{PAGA} % whatever is in step bible without periods
			
				% \begin{compactdesc}
					% \item[HALOT , PDF ] \textbf{} % Use the 1994 PDF
						% \begin{compactitem}
							% \item[1]
						% \end{compactitem}
					% \item[BDB , PDF ] \textbf{}
						% \begin{compactitem}
							% \item[1]
						% \end{compactitem}
					% \item[DCH PDF ] \textbf{}
						% \begin{compactitem}
							% \item[1]
						% \end{compactitem}
				% \end{compactdesc}
				
				% \begin{compactdesc}
					% \item[Some OT uses] \textbf{}
						% \begin{compactdesc}
							% \item[]
						% \end{compactdesc}
				% \end{compactdesc}
			
		% % --
		% \paragraph{Septuagint}
		% % --
			% \label{RESURRECTION-LXX}
		
			% %
			% \subparagraph{\texorpdfstring{\gr{sample word}}{}}
			% %
		
	% -----
	\subsubsection{New Testament} % *NT
	% -----
		\label{RESURRECTION-NT}
	
		% --
		\paragraph{KJV}
		% --
			\label{RESURRECTION-NT-KJV}
			
	
			% \begin{tabular}{ll}
				% Total occurrences in the NT & 
			% \end{tabular}
			
			\begin{compactdesc}
				\item[{\hyperref[LUKE-14-12]{Luke 14:12--14}}]
				\item[{\hyperref[JOHN-5-17]{John 5:17--47}}]
				\item[{\hyperref[JOHN-20]{John 20}}] \phantom{.}
					\begin{compactitem}
						\item Note that there are many important things happening in these verses, including what Christ says about not touching him because he hasn't ``ascended'' yet.
					\end{compactitem}
				\item[{\hyperref[ROMANS-6]{Romans 6}}]
				\item[{\hyperref[2CORINTHIANS-4-14]{2 Corinthians 4:14}}]
				\item[{\hyperref[PHILIPPIANS-3-10]{Philippians 3:10--11}}] \phantom{.}
					\begin{compactitem}
						\item The larger context is in verses 8--14, but 10 and 11 both mention the word ``resurrection.'' 
						\item Note that the words in the two verses are different---you get the expected \gr{ἀνάστασις} in 10, but then 11 is \gr{ἐξανάστασις}; see BDAG 305a.
					\end{compactitem}
			\end{compactdesc}
			
		% % --
		% \paragraph{Greek New Testament} % *GREEK
		% % --
			% \label{RESURRECTION-GREEK}
		
			% %
			% \subparagraph{\texorpdfstring{\gr{sample word}}{}}
			% %
				% \label{KATALLAGE}
			
				% \begin{compactdesc}
					% \item[BDAG , PDF ] \textbf{}
						% \begin{compactitem}
							% \item 
						% \end{compactitem}
					% \item[CGL , PDF ] \textbf{}
						% \begin{compactitem}
							% \item 
						% \end{compactitem}
					% \item[LSJ , PDF ] \textbf{}
						% \begin{compactitem}
							% \item 
						% \end{compactitem}
				% \end{compactdesc}
				
				% \begin{compactdesc}
					% \item[Some NT uses] \textbf{}
						% \begin{compactdesc}
							% \item[]
						% \end{compactdesc}
				% \end{compactdesc}
		
	% -----
	\subsubsection{Book of Mormon} % *BOM
	% -----
		\label{RESURRECTION-BOM}
	
		% \begin{tabular}{ll}
			% Total occurrences in the BOM & 
		% \end{tabular}
		
		\begin{compactdesc}
			\item[{\hyperref[2NEPHI-9-4]{2 Nephi 9:4}}] For I know that thou hast searched much, many of you, to know of things to come. Wherefore I know that ye know that our flesh must waste away and die; nevertheless in our bodies we shall see God.
				\begin{compactitem}
					\item Not ``in our flesh we shall see God,'' but ``in our bodies''; not ``with our bodies'' here. It's kind of an unusual comment. It makes me think of \hyperref[ALMA-11-45]{Alma 11:45}, where ``the whole'' (body) becomes ``spiritual and immortal.'' \hyperref[2NEPHI-9-5]{2 Nephi 9:5} gives a clue about the meaning of ``in the body,'' because it talks about Christ showing himself to people at Jerusalem in that way.
				\end{compactitem}
			\item[{\hyperref[ALMA-11-45]{Alma 11:45}}]
				\begin{compactitem}
					\item And other verses around this one, but I think this one may be the most important.
				\end{compactitem}
			\item[{\hyperref[HELAMAN-14-15]{Helaman 14:15--19}}]
			\item[{\hyperref[3NEPHI-26-4]{3 Nephi 26:4--5}}]
			\item[{\hyperref[3NEPHI-28-36]{3 Nephi 28:36--40}}] \phantom{.}
				\begin{compactitem}
					\item This passage is super crucial.
				\end{compactitem}
			\item[{\hyperref[MORMON-7-5]{Mormon 7:5--7}}]
			\item[{\hyperref[MORMON-9-12]{Mormon 9:12--14}}]
		\end{compactdesc}
		
	% -----
	\subsubsection{Doctrine and Covenants} % *DC
	% -----
		\label{RESURRECTION-DC}
	
		% \begin{tabular}{ll}
			% Total occurrences in the DC & 
		% \end{tabular}
		
		\begin{compactdesc}
			\item[{\hyperref[DC29-11]{DC 29:11--13}}] 
			\item[{\hyperref[DC45-45]{DC 45:45--46}}]
			\item[{\hyperref[DC45-54]{DC 45:54}}]
			\item[{\hyperref[DC52-44]{DC 52:44}}]
			\item[{\hyperref[DC63-17]{DC 63:17--18}}]
			\item[{\hyperref[DC63-49]{DC 63:49--52}}]
			\item[{\hyperref[DC76-39]{DC 76:39}}]
				\begin{compactitem}
					\item There are other verses around this area which could be relevant, like 37.
					\item Beginning in verse 50, the discussion of the resurrection turns to that of the ``just,'' and continues from there.
					\item There's a lot about the resurrection in this section.
					\item See 64--65 and notes there too.
					\item And 85.
				\end{compactitem}
			\item[{\hyperref[DC88-14]{DC 88:14--35}}]
			\item[{\hyperref[DC88-95]{DC 88:95--116}}]
			\item[{\hyperref[DC138-51]{DC 138:51}}]
		\end{compactdesc}
		
	% % -----
	% \subsubsection{Pearl of Great Price} % *PGP
	% % -----
		% \label{RESURRECTION-PGP}
	
		% \begin{tabular}{ll}
			% Total occurrences in the PGP & 
		% \end{tabular}
		
		% \begin{compactdesc}
			% \item[]
		% \end{compactdesc}
		
% % ----------------------
% \subsection{Key Phrases} % *KEYPHRASES *PHRASES
% % ----------------------
	% \label{RESURRECTION-PHRASES}

	% \begin{compactdesc}
		% \item[] \textbf{\#times} | list
			% % \begin{compactdesc}
				% % \item[Bible anchor verses] \phantom{.}
					% % \begin{compactdesc}
						% % \item[Romans 5:5] \gr{}
					% % \end{compactdesc}
				% % \item[Commentary] ---
			% % \end{compactdesc}
	% \end{compactdesc}
	
% % ----------------------
% \subsection{Bible Dictionaries} % *DICTIONARIES
% % ----------------------
	% \label{RESURRECTION-BD}

% ----------------------
\subsection{Restoration Usage} % *RESTORATION
% ----------------------
	\label{RESURRECTION-RESTORATION}

	% -----
	\subsubsection{Joseph Smith} % *JS
	% -----
		\label{RESURRECTION-JS}
	
		\begin{compactdesc}
			\item[{\hyperref[EMS-1843-JS-23-MAR-1843-JK]{23 March 1843}}] \phantom{.}
				\begin{compactitem}
					\item This is a blessing JS gave to Joseph Kingbury, in which he says a lot of things about the resurrection.
				\end{compactitem}
			\item[{\hyperlink{EMS-1843-JS-6-8-APR-1843-WR-i}{6--8 April 1843 WR-i}}] <P.> Joseph said to complete the subject of Bro Pratts. I thought it a glorious subject with one \sout{additional idea} <addition> their is no fundamental principle belonging to a human System that \sout{n}ever goes into another--- <in this world. or the world to come.---> the principle of Mr Pratt was correct. I care not what the theories of \sout{men} <man> are--- we have the testimony that God will raise us up \& he has power to do it.--- If any one supposes--- that any part of our bodies. that is the fundame[n]tal parts thereof, ever goes into another body \sout{t}he is mistaken
				\begin{compactitem}
					\item Orson Pratt had been speaking about the resurrection; see his remarks \hyperref[EMS-1843-OP-6-8-APR-1843-WR]{here}.
				\end{compactitem}
		\end{compactdesc}
		
	% % -----
	% \subsubsection{Temple Ordinances}
	% % -----
		% \label{RESURRECTION-TEMPLE}
	
		% \begin{compactdesc}
			% \item[{\hyperref[KEY]{Date/Source}}]
		% \end{compactdesc}
	
	% -----
	\subsubsection{Brigham Young} % *BY
	% -----
		\label{RESURRECTION-BY}
	
		\begin{compactdesc}
			\item[6 October 1867] no man can explain the resurrection until they receive the keys of the resurrection no more than John Wesley could explain the gospel and building up of Kingdom of God in the latter day without the priesthood and you know he could tell nothing about it [audiences?] to be begin with we talk of scriptures for what other men may understand things by the revelations of the Lord Jesus but they cannot explain to unbelievers all the things of God are understood by the spirit of God without it we know nothing of the resurrection Paul in his writings to his brethren I suppose labored as faithfully to explain principle of the resurrection as a man possibly could do and he utterly failed to give true ideas to his brethren that is in some things he could make a great many statements that are true but still he had not the keys of the resurrection neither could he administer and explain to others \textit{(See this source online \href{https://drive.google.com/file/d/1wV-mRM932G6giTCQENRd8CTtRAMDGdxa/view?usp=sharing}{here}.)}
		\end{compactdesc}
	
% ----------------------
\subsection{Commentary and Studies} % *COMMENTARY *STUDIES
% ----------------------
	\label{RESURRECTION-COMMENTARY}

	% -----
	\subsubsection{Key Points}
	% -----
		\label{RESURRECTION-KEYS}
	
		\begin{compactitem}
			\item The resurrection is \textit{universal}, except the sons of perdition, apparently.
				\begin{compactitem}
					\item \hyperref[ALMA-11-42]{Alma 11:42--44}, \hyperref[DC76-36]{DC 76:36--39}
				\end{compactitem}
		\end{compactitem}
		
	% -----
	\subsubsection{Bodies being ``spiritual'' after the resurrection}
	% -----
		\label{RESURRECTION-spiritual}
		
		\begin{compactdesc}
			\item[Augustine, DDC I:48--49] Some say that they would prefer not to have a body at all, but they are mistaken. For what they hate is not their body, but its imperfections and its dead weight. 49. What they want is not to have no body at all, but to have one free from corruption and totally responsive; they think that if the body were such a thing it would not be a body, because they consider such a thing to be a soul. When they seem to persecute their own body by a kind of repression, and by hardships, their aim (if they are doing it rightly) is not to have no body at all but to have one that is subservient and ready for necessary tasks. (See also up through 53.)
			\item[{\hyperref[EMS-1865-BY-8-OCT-1865]{BY, 8 October 1865}}] I don't believe that the Lord is a sinful man neither do you any of you I believe he is a glorified man an exalted man a spiritual man I will give you a comparison I will not stray from the belief of the Christian world you go to them and ask [if] they believe in the resurrection yes they tell you pretty much all the sects they believe that the body will be resurrected immortal spirit and be crowned those that have lived their religion no matter what it is Catholic Church [of] England or any of the religions yes if we live according to our creed we shall receive a resurrection when we see this resurrected person according to their faith and according to ours we see a spiritual body it is sown in weakness raised in power sown a natural body raised a spiritual body now this you and I believe the whole Christian world with few exceptions believe the same and those that receive a resurrection will be raised immortal spiritual incorruptible this is correct we believe it because it is true 
		\end{compactdesc}
		
	% -----
	\subsubsection{Stories to which we can apply resurrection thinking}
	% -----
		\label{RESURRECTION-stories}
		
		\begin{compactitem}
			\item \hyperref[ALMA-22-22]{Alma 22:22} (and compare when Abish and the queen do it a few chapter earlier)
		\end{compactitem}